\section{Effectiveness criteria for a descent datum}
\label{VIII.7}

Consider as usual a morphism of preschemes
$$
g\colon S'\to S
$$
and an $S'$-prescheme~$X'$. In accordance with the
general facts~VII,~9\footnote{does not exist; see expos\'e 212 of the
S\'eminaire Bourbaki (1962)},
giving a descent datum on~$X'$, relative
to~$g$, is equivalent to giving an equivalence
pair \cite{VIII.3}:
$$
\xymatrix@C=.5cm{ q_{1},q_{2}\colon X''\ar@<2pt>[r]\ar@<-2pt>[r] & X'}
$$
such that the structural morphism~$X'\to S'$ is compatible with this
pair and the equivalence pair
$$
\xymatrix@C=.5cm{p_{1},p_{2}\colon
S''=S'\times_{S}S'\ar@<2pt>[r]\ar@<-2pt>[r] & S'}
$$
defined by~$g$, and such that the two squares (or one of the two,
which amounts to the same by reason of symmetry) extracted from the
corresponding diagram
$$
\xymatrix{X'\ar[d] & X''\ar@<2pt>[l]\ar@<-2pt>[l]\ar[d]\\ S' &
S''\ar@<2pt>[l]\ar@<-2pt>[l]}
$$
(using either~$p_{1},q_{1}$, or~$p_{2},q_{2}$), are
\emph{cartesian}. A solution of the descent problem
posed by this descent datum, \ie an object~$X$ over~$S$
endowed with an isomorphism~$X\times_{S}S'\isomfrom X'$ compatible with the
descent data, is equivalent to giving a
\emph{cartesian} square
$$
\xymatrix{X\ar[d] & X'\ar[l]_{h}\ar[d]\\ S & S'\ar[l]_{g}}
$$
satisfying~$hq_{1}=hq_{2}$.

Since
% original p. 220
the set of faithfully flat and quasi-compact morphisms is
stable under change of base, and since a faithfully flat
and quasi-compact morphism is an effective epimorphism by virtue
of~\Ref{VIII.5.3}, it follows from the general theory
\cite{VIII.D}:

\begin{proposition}
\label{VIII.7.1}
Suppose~$g\colon S'\to S$ faithfully flat and quasi-compact. For
a descent datum on~$X'$ relative to~$g$ to be
effective, it is necessary and sufficient that the equivalence
relation~$R=(q_{1},q_{2})$ which it defines be effective
(\ie the quotient~$X'/R$ exists and $X''$ becomes the fibered
square of~$X'$ over~$X'/R$), and that the canonical morphism~$X'\to
X'/R$ be faithfully flat and quasi-compact.
\end{proposition}

Thus the question of the effectiveness of a descent datum
is a special case of the question of effectiveness of an equivalence
graph, and various effectiveness criteria given
in this no.\ can be obtained in this way. Nevertheless, in the context
of descent one has at one's disposal
theorem~\Ref{VIII.2.1}, which implies that \emph{if $X'$ is
affine over~$S'$, every descent datum on~$X'$ relative
to~$g$ is effective}, a statement which has no analogue for
passage to the quotient by a general flat equivalence
graph. All the effectiveness criteria that we
give here can also be considered as
deduced from the preceding one.

Let~$U'$ be a subprescheme of~$X'$ (or more
generally a subobject of~$X'$ in the
category~$\Sch$); one says that~$U'$ is \emph{stable under the
descent datum} on~$X'$ if one can find on~$U'$ a
descent datum relative to~$g$ such that
the immersion~$U'\to X'$ is compatible with the descent
data. This also means that the inverse images of~$U'$
in~$X''$ under~$q_{1}$ and~$q_{2}$ are the same (or also, as
one says, that~$U'$ is \emph{stable under the equivalence
relation}~$R$), and of course the descent datum in
question on~$U'$ is then unique, and is called the \emph{descent
datum induced}
by that of~$X'$. This being said:

\begin{proposition}
\label{VIII.7.2}
Let~$(X_{i}')$ be a covering of~$X'$ by open sets~$X_{i}'$
stable
% original p. 221
under the descent datum. For the descent datum
on~$X'$ to be effective, it is necessary and sufficient that the
same hold for the induced descent data on the~$X_{i}'$.
\end{proposition}

This is an easy consequence of~\Ref{VIII.7.1} for instance,
and the details of the proof are left to the reader.

\begin{corollary}
\label{VIII.7.3}
Let~$(S_{i})$ be an open covering of~$S$, and for every~$i$
let~$S_{i}'$ and~$X_{i}'$ be deduced from~$S'$
and
$X'$ by the change of base~$S_{i}\to S$. For the descent
datum on~$X'$ to be effective, it is necessary and sufficient that, for
every~$i$, the descent datum on~$X_{i}'$ relative
to~$g_{i}\colon S_{i}'\to S_{i}$ be effective.
\end{corollary}

This criterion practically always reduces us to the case where~$S$
is affine. In the case where~$S'$ is also affine, which
is the most frequent case in applications, one has:

\begin{corollary}
\label{VIII.7.4}
Suppose~$S$ and~$S'$ affine. For the descent datum
on~$X'$ to be effective, it is necessary and sufficient that~$X'$ be
a union of open sets~$X_{i}'$ that are affine and stable under the
descent datum.
\end{corollary}

The sufficiency comes from~\Ref{VIII.7.2} and from the fact that if~$X_{i}'$
is affine, it is affine over~$S'$ and one can
apply~\Ref{VIII.2.1}. For the necessity, one notes that
if~$X'$ comes from~$X$, and if~$X$ is covered by affine open
sets~$X_{i}$, then the~$X_{i}'=X_{i}\times_{S}S'$ are affine open
sets stable under the descent data and covering~$X'$.

\begin{corollary}
\label{VIII.7.5}
Let~$g\colon S'\to S$ be a faithfully flat, quasi-compact
and \emph{radicial} morphism. Then~$g$ is an
\emph{effective} descent morphism, \ie for every~$X'$ over~$S'$, every
descent datum on~$X'$ relative to~$g\colon S'\to S$ is effective.
\end{corollary}

Indeed, by virtue of~\Ref{VIII.7.3} one may assume~$S$ affine; hence,
as~$S'$ is radicial over~$S$ and hence separated, $S'$ is
separated. Moreover, for every~$x'\in X'$, the
fiber~$R(x')=q_{2}(q_{1}^{-1}(x'))$ of the set-theoretic equivalence
relation defined
% original p. 222
by the equivalence relation~$R$ is reduced to a point,
because, $g$ being radicial, the same holds for~$p_{1},p_{2}$, which
are deduced from it by the change of base~$S'\to S$, hence also
for~$q_{1},q_{2}$, which are deduced from the preceding ones by the
change of base~$X'\to S''$. Hence \emph{every open set of~$X'$ is
stable} under the descent datum. Let us then cover~$X'$ by
affine open sets~$X_{i}'$; they are affine over~$S$ since~$S'$ is
separated, hence the induced descent datum is effective
by~\Ref{VIII.2.1}. One then concludes by~\Ref{VIII.7.2}.

One will note that~\Ref{VIII.7.5} gives the only known case of an
effective descent morphism in the category of preschemes, and
it is probably indeed the only case, even when restricting oneself to
noetherian schemes, or to schemes of finite type over a
field.

When one assumes~$S$ locally noetherian and~$S'$ of finite type
over~$S$, the statement~\Ref{VIII.7.5} is also a special case
of the following one (which generalizes Weil's Galois descent and
Cartier's inseparable descent):

\begin{corollary}
\label{VIII.7.6}
Let~$g\colon S'\to S$ be a finite locally free morphism
(\ie defined by an Algebra over~$S$ which is a
locally free module of finite type) and surjective (hence~$g$ is
faithfully flat and quasi-compact, hence a descent
morphism). Let~$X'$ be an~$S'$-prescheme endowed with a descent
datum. For this datum to be effective, it is necessary and
sufficient that for every~$x'\in X'$, the
fiber~$R(x')=q_{2}(q_{1}^{-1}(x'))$ be contained in an affine open
set (a condition automatically satisfied if~$X'$ is
quasi-projective over~$S'$).
\end{corollary}

The remark in parentheses comes from the fact that if~$s$ is the
point of~$S$ below~$x'$, then $R(x')$ is finite and contained
in the
fiber~$X'_{s}$;
on the other hand, since~$X'$ is quasi-projective over~$S'$ and~$S'$ is finite
over~$S$, $X'$ is quasi-projective over~$S$, which implies that a
fiber
of~$X'$ over~$S$
is contained in an affine open set.

Since every finite subset of an affine scheme admits a fundamental
system of affine neighborhoods, one sees that one does not lose
the hypothesis by restricting oneself to the part over an affine open set
of~$S$, which by virtue of~\Ref{VIII.7.3} reduces us to the case
where~$S$ is affine. By virtue of~\Ref{VIII.7.4}, we are reduced
to showing
that~$x'$
is contained
% original p. 223
in an affine open set \emph{stable} under the descent
datum. Indeed, let~$U$ be an affine open set containing~$R(x')$; then
the saturation
$$
R(X'-U)=q_{2}(q_{1}^{-1}(X'-U))
$$
does not meet~$R(x')$; on the other hand, since~$q_{2}$ is finite (because~$g$,
hence~$p_{2}$, is) and therefore closed, the right-hand side is a
closed subset of~$X'$. Let~$U'$ be its complement in~$X'$;
it is therefore a \emph{saturated} open set and one has
$$
R(x')\subset U'\subset U
$$
with~$U$ affine, but~$U'$ not affine a priori. Since a finite
subset~$R(x')$ in an affine scheme~$U$ has a fundamental
system of affine neighborhoods of the form~$U_{f}$, one sees,
replacing~$f$ by its restriction to~$U'$, that there exists a
section~$f$ of~$\mathcal{O}_{U}$ such that:
$$
R(x')\subset U_{f}',\quad U_{f}'\textrm{ is affine}.
$$
Then let~$U''=q_{1}^{-1}(U')=q_{2}^{-1}(U')$, let us again denote
by~$q_{1},q_{2}$ the induced morphisms~$U''\to U'$, and
consider
$$
f'=\Norm_{q_{2}}(q_{1}^*(f))\quoi,
$$
where~$\Norm_{q_{2}}$ denotes the \emph{norm} relative to the
finite locally free morphism~$q_{2}\colon U''\to U'$. The
compatibility of the formation of the norm with change
of base easily implies that~$f'$ is an \emph{invariant} section:
$$
q_{1}^*(f')=q_{2}^*(f')
$$
which implies that~$U_{f'}'$ is a saturated open set of~$U'$. More
precisely, in fact, denoting by~$Z(f')$
the set of zeros of a section~$f'$, one finds, by virtue of the
properties of norms:
$$
Z(f')=q_{2}(Z(q_{1}^*(f)))=q_{2}(q_{1}^*(Z(f)))=
R(U'-U_{f}')\quoi,
$$
which
% original p. 224
implies that~$U_{f'}'=U'-Z(f')$ is saturated, contains~$R(x')$,
and is contained in~$U_{f}'$. Since the latter is affine, it follows
that~$U_{f'}'$ is affine too (since it is equal to~$(U_{f}')_{f''}$,
with~$f''=f'|U_{f'}'$). It is therefore a saturated affine open set
containing~$R(x')$, hence~$x'$, which completes the proof.

One will note that this reasoning applies whenever one has an
equivalence relation (or even only a pre-equivalence relation,
\cf \cite{VIII.3}) in a prescheme~$X'$, finite
and locally free, and moreover~\Ref{VIII.7.6} is also a special
case of the analogous result for finite and locally free
pre-equivalences, \cf \loccit The same remark
applies to~\Ref{VIII.7.7} below.

One can also, once the existence of a saturated quasi-affine open
set~$U'$ containing~$x'$ has been obtained, appeal to~\Ref{VIII.7.9}
and~\Ref{VIII.7.2}, which avoids recourse to norms.

Let us note moreover that under the conditions of~\Ref{VIII.7.6}, if the
descent datum on~$X'$ is effective, $X'$ coming from~$X$
over~$S$, then the morphism~$X'\to X$ is finite, locally free and
surjective, since it is deduced from~$g$ by the change of base~$X\to
S$. It follows (EGA~II 6.6.4) that if~$X'$ is quasi-projective
over~$S'$, hence over~$S$, then~$X$ is quasi-projective over~$S$ (a
relatively ample invertible sheaf on~$X$ being obtained by
taking the \emph{norm} of an invertible sheaf on~$X'$ relatively
ample over~$S$, or over~$S'$, which amounts to the same). One thus
obtains:

\begin{corollary}
\label{VIII.7.7}
A finite locally free and surjective morphism~$g\colon S'\to S$ is
an effective descent morphism for the fibered category
of preschemes quasi-projective over other ones, \ie for
every~$X'$ quasi-projective over~$S'$, every descent datum
on~$X'$ relative to~$g$ is effective, and the
descended $S$-prescheme~$X$ is quasi-projective over~$S$.
\end{corollary}

\begin{proposition}
\label{VIII.7.8}
Let~$g\colon S'\to S$ be a faithfully flat and
quasi-compact morphism. Then~$g$ is an effective descent morphism for
the fibered category of preschemes~$Z$ quasi-compact
over a prescheme~$T$, endowed with an invertible sheaf ample
% original p. 225
relative to~$T$. In particular, for every
prescheme~$X'$ over~$S'$ endowed with a descent datum
relative to~$g\colon S'\to S$, and every invertible
sheaf~$\mathcal{L}'$ on~$X'$ ample relative to~$S'$, also
endowed with a descent datum relative to the one
given on~$X'$ (\ie endowed with an isomorphism
of~$q_{1}^*(\mathcal{L}')$ with~$q_{2}^*(\mathcal{L}')$
satisfying the usual transitivity condition), the
descent datum on~$X'$ is effective, and the invertible
sheaf~$\mathcal{L}$ on the descended prescheme~$X$,
deduced from~$\mathcal{L}'$ by descent, is ample relative
to~$S$.
\end{proposition}

The proof is quite analogous to that of~\Ref{VIII.5.8},
noting that on the quasi-coherent graded
Algebra~$\mathcal{S}'$ over~$S'$ defined
by~$\mathcal{L}'$ there is a descent datum, making it possible to
construct a quasi-coherent graded
Algebra~$\mathcal{S}$ over~$S$ thanks
to~\Ref{VIII.1.1}, whence a~$P=\Proj(\mathcal{S})$ over~$S$ such
that~$P'=\Proj(\mathcal{S}')$ is identified, together with its descent
datum, with~$P\times_{S}S'$. Since by hypothesis $X'$
is identified with an open set of~$P'$, necessarily stable under the
descent datum on~$P'$, the descent datum on~$X'$
is also effective, and one obtains the descended
prescheme as an open set in~$P$. The details are left to the
reader. --- In particular, taking~$\mathcal{L}'=\mathcal{O}_{X'}$, one
finds:

\begin{corollary}
\label{VIII.7.9}
Let~$g\colon S'\to S$ be a faithfully flat and
quasi-compact morphism, and let~$X'$ be a \emph{quasi-affine}
prescheme over~$S'$; then every descent datum on~$X'$
relative to~$g$ is effective, and the descended
prescheme~$X$ is quasi-affine over~$S$.
\end{corollary}

By virtue of~\Ref{VIII.6.2}, this result applies in particular
if~$S'$ is locally noetherian and~$X'$ is quasi-finite and
separated over~$S'$, more generally if~$S'$ is
arbitrary and~$X'$ is of finite presentation, quasi-finite and
separated over~$S'$ (\cf \Ref{VIII.6.6}).

\begin{remarks}
\label{VIII.7.10}
The results given in this no.\ exhaust the effectiveness
criteria currently known, and probably
even the useful criteria that exist\footnote{This opinion
has turned out to be partially mistaken, see for example J.P\ptbl \textsc{Murre},
S\'em. Bourbaki 294 (Appendix), May 1965, and special
results (notably of \textsc{N\'eron} and \textsc{Raynaud}). For the descent of
group schemes; \cf M\ptbl Raynaud, thesis (1968).}. One will note
the following counterexamples in support of this assertion:
\begin{enumerate}
\item[(i)]
% original p. 226
If~$S$ is the spectrum of a field, and~$S'$ is the spectrum of a
Galois quadratic extension, one can find an~$X'$ over~$S'$,
proper and smooth over~$S'$, of dimension 3, endowed with a
descent datum which is not effective (Serre).
\item[(ii)] One can find an~$S$ spectrum of a regular local ring
of dimension 3 (if one wishes, the local ring of an
algebraic scheme over a field of given
characteristic), a~$T$ principal covering of~$S$ with
group~$\ZZ/2\ZZ$, such that, if~$t$ denotes one of the points of~$T$
above the closed point~$s$ of~$S$, and $S'=T-s$, one can
find an~$X'$ \emph{projective} over~$S'$, regular, endowed with a
descent datum relative to~$g\colon S'\to S$, this
descent datum not being effective.
\end{enumerate}

For these constructions one uses Hironaka's example of
non-projective varieties. For~(i), it suffices to use the
fact that one can find over~$k$ a proper and
smooth scheme~$X_{0}$ of dimension 3, on which~$G=\ZZ/2\ZZ$ acts without
inertia, and in which there exist two points~$a,b$ rational
over~$S$,
congruent under~$G$, which are not contained in an affine open set. One
then puts~$X'=X_{0}\times_{k}k'$, and one lets~$G$ act on~$X'$
by means of the actions of~$G$ on the two factors, which
gives a descent datum on~$X'$ relative to~$g\colon
\Spec(k')\to \Spec(k)$. Above~$a$ \resp $b$, there is exactly
one point~$a'$ \resp $b'$ (with quadratic residue field
extension), and~$a'$ and~$b'$ are congruent under~$G$, since~$X'\to
X_{0}$ is compatible with the actions of~$G$. Then~$a'$
and~$b'$ cannot be contained in an affine open set, say~$U'$,
for then~$U=X_{0}-\Im(X'-U')$ would be an open set of~$X_{0}$
containing~$(a,b)$ and whose inverse image in~$X'$ would be contained
in~$U'$, hence quasi-affine, hence~$U$ would be quasi-affine, and
consequently~$(a,b)$ would have an affine neighborhood in~$U$.

For~(ii), one uses the fact that in Hironaka's example, $X_{0}$
is obtained as a prescheme proper over
a projective $k$-scheme~$Y$, smooth over~$k$ (the
morphism~$f\colon X_{0}\to Y$ being moreover birational, but
that does not matter), the group~$G$ also acting on~$Y$
compatibly with its actions on~$X_{0}$; finally,
putting~$S'=Y-f(b)$, $X'=X_{0}|S'$, $X'$ is projective
over~$S'$. Then~$X_{0}$ is endowed with a natural descent
datum relative to the canonical
% original p. 227
morphism~$Y\to S=Y/G$, thanks to the actions of~$G$
on~$X_{0}$ compatible with its actions on~$Y$. This
descent datum is not effective, since~$(a,b)$ is not
contained in an affine open set. The induced descent datum
on~$X'$ relative to~$g\colon S'\to S$ is then not
effective, as one easily verifies.
\end{remarks}

\begin{thebibliography}{D}{VIII.8}

\bibitem[D]{VIII.D} J\ptbl \textsc{Giraud}, M\'ethode de la descente, M\'emoire \no 2
de la Soc.\ Math.\ Fr., 1964.

\bibitem[1]{VIII.1} A\ptbl \textsc{Grothendieck}, S\'eminaire Bourbaki:
G\'eom\'etrie formelle et G\'eom\'etrie alg\'ebrique, May
1959, \No 182.

\bibitem[2]{VIII.2} A\ptbl \textsc{Grothendieck}, S\'eminaire Bourbaki: Technique de
descente et Th\'eor\`emes d'existence~I, December 1959,
\No 190.

\bibitem[3]{VIII.3} A\ptbl \textsc{Grothendieck}, S\'eminaire Bourbaki: Technique de
descente et Th\'eor\`emes d'existence~III, February 1961,
\No 212.
\end{thebibliography}
